% ===================================================================
%  HOTO Na 7 — Ƙauye da hunturu. --- Gidan gona da bazara.
%
%  Traduction, faite en 2026, en haoussa d'aujourd'hui, sur l'ido de
%  texto/io. Regles de la colonne : voir 10-hoto-01.tex. Deux scenes,
%  deux numerotations. Ouverture de la DEUXIEME SERIE : « JERI NA
%  BIYU ». Un bloc marque « ¶2 » par ordo.py, t07-c1-07-1 : une seule
%  cle, deux alineas — ne pas en ouvrir une seconde.
%
%  LA BASSE-COUR EST LE TROISIEME LEXIQUE QUE LE HAOUSSA NE DOIT A
%  PERSONNE, apres le corps du tableau 2 et la cuisine du tableau 6 :
%
%    doki       le cheval     godiya    la jument
%    saniya     la vache      bijimi    le taureau
%    tunkiya    la brebis     rago      le belier
%    akuya      la chevre     alade     le porc
%    kaza       la poule      zakara    le coq
%    agwagwa    le canard     talotalo  le dindon
%    tantabara  le pigeon     zomo      le lapin
%    gona       le champ      noma      labourer
%    taki       le fumier     rijiya    le puits
%    kwatarniya l'auge        garke     le troupeau
%    makiyayi   le berger     manomi    le fermier
%    nono       le lait       ƙwai      l'oeuf
%
%  Pas un emprunt. Ces trois lexiques-la — ce qu'on a dans son corps,
%  ce qu'on cuit et ce qu'on eleve — sont les trois seuls du livret
%  ou la colonne n'ait rien eu a chercher.
%
%  ET « keke » PARAIT POUR LA QUATRIEME FOIS, CE QUI MESURE
%  L'ANNONCE DU TABLEAU 2. On avait ecrit la que keke etait le mot de
%  la machine a manivelle et non un emprunt fait pour la bicyclette.
%  Le releve peut maintenant le compter : keke la bicyclette (tableau
%  2), keken doki la voiture attelee (tableau 5), keken kaɗi le rouet
%  et keke la roue du moulin (17) ici. Quatre objets, un seul mot, et
%  la seule chose qu'ils aient en commun est de TOURNER SUR UN AXE.
%
%  SIXIEME REFUS DE LA COLONNE, ET C'EST LE PLUS SERRE DE TOUS : LES
%  MARRONS (42). « gyaɗa » est l'arachide grillee, et la scene gravee
%  est mot pour mot celle qu'on voit a Kano des que l'harmattan
%  tombe : une vieille femme qui grelotte sous un pagne trop mince,
%  un brasero, un cornet de papier, des enfants qui se rechauffent les
%  mains avec une poignee brulante achetee deux sous. Meme scene,
%  meme geste, meme saison, meme commerce — et ce n'est pas la meme
%  graine. On ecrit « kastanya ». La ressemblance parfaite d'un usage
%  ne fait pas l'identite d'une chose.
%
%  LE CYPRES (22) EST LE TROISIEME ARBRE QUE LE SAHEL N'A PAS, ET LA
%  PRUNE (15) LE TROISIEME FRUIT QU'IL N'A PAS : la regle posee au
%  tableau 3 tient donc sur quatre objets de plus, et dans les deux
%  sens. Le peuplier, le platane et le cypres sont DECRITS —
%  « dogon icen makabarta », le grand arbre du cimetiere. La peche,
%  l'abricot et la prune sont EMPRUNTES et refaits au p — fis,
%  abirkot, furam. Le fruit arrive dans une caisse et l'arbre
%  n'arrive jamais.
%
%  LE SABOT DU CHEVAL SE DIT « kofato », ET LA REGLE NE CRIE PAS. Le
%  controle releve « kofa » nu, qui n'est rien en haoussa la ou
%  « ƙofa » est la porte ; « kofato » commence par les memes quatre
%  lettres et n'a rien a voir. La regle est ecrite avec des frontieres
%  de mot des deux cotes, et c'est pour ce genre de mot-la qu'elles y
%  sont. Un controle se juge autant sur ce qu'il laisse passer.
%
%  DEUX MOTS DE METIER QUI VALENT UNE PHRASE : « masauki »
%  l'auberge (2), le lieu ou l'on est loge, encore un ma- ; et
%  « ƴan sanda » la police, les FILS DU BATON, dont le tableau tire
%  son ancien gendarme (27). Et le barbier (8) est « wanzami », qui
%  est un vrai metier haoussa, hereditaire, celui qui rase, coupe,
%  circoncit et saigne : on l'accepte sans hesiter, parce que
%  l'homme de la planche fait exactement ce que fait le wanzami.
%
%  UNE FAUTE D'ORDRE, LA TROISIEME, ET ELLE N'EST PAS DE LA MEME
%  ESPECE QUE LES DEUX PREMIERES. Au bloc t07-c2-13-1 l'ido ecrit
%  « La tota lakteyo (59) mantenesas tre nete da la farm-servisto
%  (65) » : une PASSIVE, qui pose le lieu en tete et rejette l'agent.
%  La colonne l'avait rendue par une active — le valet, puis la
%  laiterie —, et l'ordre s'etait retourne. Les deux fautes du
%  tableau 3 et du tableau 6 etaient des choix de tour, que rien
%  n'imposait ; celle-ci vient d'une PRESSION de la langue, qui
%  prefere l'actif. Le remede n'est pas de forcer un passif que le
%  haoussa n'a pas, c'est l'impersonnel : « ana kula da... ta wurin
%  ma’aikacin gona », on entretient... par le valet. Trois fautes en
%  sept tableaux, deux de placement et une de voix.
%
%  UN \VUgras{} VIDE, LE PREMIER DE LA COLONNE, ET C'EST L'ESPECE
%  ORDINAIRE : au bloc t07-c2-10-1 le renvoi du chien (55) suivait un
%  groupe gras deja ferme — « \VUgras{Karensa} \VUgras{} » — et le
%  second n'avait rien a mettre en gras. Les colonnes indonesienne et
%  javanaise en avaient fait quatre a elles deux, toujours au meme
%  endroit : un renvoi qui suit un terme dont le gras est clos.
%  kolonoj.py l'a releve au premier passage.
%
%  ENFIN L'AUBERGE S'APPELLE « L'OURS BLANC », « Farin Beyar », et
%  c'est l'emprunt du tableau 2 qui revient en enseigne. Un mot pris
%  pour une bete qu'aucun Haoussa n'a vue sert ici a peindre un
%  panneau qui grince au vent du nord.
% ===================================================================

%%K t07-noto-1 noto
\VUnotes{143.2mm}{%
(*) Domin cikakken bayani game da abinci, duba hoto na 6 da na 13.%
}

%%K t07-apar-1 apar
\VUsaut{4.0mm}
\VUcentre{13.2pt}{90}{JERI NA BIYU}
\VUsaut{3.0mm}
\VUfilet{16mm}
\VUsaut{5.6mm}
\VUtitre{15.0pt}{60}{HOTO Na 7}
\VUsaut{3.0mm}

%%K t07-tit-1 sub
\VUcentre{12.0pt}{0}{\textit{Fage na Farko.}}
\VUsaut{3.0mm}

%%K t07-tit-2 sub
\VUpk{10.2pt}{40}{Ƙauye a lokacin hunturu.}
\VUsaut{4.0mm}

%%K t07-c1-01-1 p
1. --- A hutun sabuwar shekara, na raka mahaifina zuwa Novurbo,
ƙaramin ƙauye inda yake buƙatar sasanta waɗansu sha’anin kasuwanci.

%%K t07-c1-02-1 p
2. --- Mun yi amfani da \VUgras{babbar karusa} \textsuperscript{(11)}
wadda take tafiya tsakanin gari da wannan ƙauye ; amma tun da ana sanyi
ƙwarai, wannan tafiya cikin abin hawa marar daɗi bai yi daɗi sosai ba.

%%K t07-c1-03-1 p
3. --- Shi ya sa muka ji daɗi na gaskiya lokacin da muka ga \VUgras{mai
tuƙi} ya tsayar da abin hawansa a kan filin, gaban \VUgras{masauki}
\textsuperscript{(2)} na \textit{Farin Beyar}. \VUgras{Mai masauki}
\textsuperscript{(4)} yana jiranmu a bakin ƙofar gidansa, yana shan
\VUgras{bututun tabarsa} cikin nutsuwa.

%%K t07-c1-03-2 p
Ya gayyace mu mu shiga, ni kuwa ban jira a roƙe ni ba sai na zauna
gaban wutar rassa da take ci a murhun bango. A lokacin na yi wa
\VUgras{iskar arewa} mai kaifi ba’a ƙwarai, wadda take sa tsohuwar
\VUgras{alamar shago} \textsuperscript{(3)} ta yi ƙara kuma take kwashe
\VUgras{hayaƙin} \textsuperscript{(1)} da yake fita daga bututun cikin
baƙaƙen guguwa.

%%K t07-c1-04-1 p
4. --- Lokacin abincin rana ne. Ba tare da ƙara jira ba, mun ci da
kyau abinci mai sauƙi amma mai daɗi da aka yi mana (*).

%%K t07-tit-3 sub
\VUpk{10.2pt}{40}{Waɗansu ziyara.}
\VUsaut{3.0mm}

%%K t07-c1-05-1 p
5. --- Bayan abincin rana, na raka mahaifina, wanda yake mai sayar da
inshora, zuwa wurin abokan cinikinsa. Da farko mun ziyarci
\VUgras{maƙeri} \textsuperscript{(5)}, wanda yake zaune kusa da masauki.
Yana cikin sa wa doki ƙarfe a ƙafa. Na yaba wa ƙarfin abokinsa wanda
ya riƙe \VUgras{kofaton} dokin da ƙarfi a kan afaronsa na fata, alhali
maƙerin yana kafa \VUgras{ƙarfen ƙafa} \textsuperscript{(6)} da
bugu masu ƙarfi.

%%K t07-c1-06-1 p
6. --- Sa’an nan mun je wurin \VUgras{mai gyaran takalma}
\textsuperscript{(7)} na ƙauyen, tsoho Iakobus. Ko da yake yana da
\VUgras{babban takalmi} mai kyau ƙwarai a matsayin alama, ya gamsu da
sabunta ƙasan wasu tsofaffin takalma da suka huda.

%%K t07-c1-06-2 p
Tsohon ya gaya mana yana dariya : « A gari, ni « mai yin takalma » ne
kuma ban da abokan ciniki. A nan, ni « mai gyaran takalma » ne, aiki
kuwa yana da yawa. »

%%K t07-c1-07-1 p
7. --- Ya yi mana magana a kan maƙwabcinsa \VUgras{wanzami}
\textsuperscript{(8)}, wanda abokan cinikinsa ba su gamsu da shi ba.
\VUgras{Askarsa} kullum sun karye, \VUgras{almakashinsa} kuwa ba sa
yanka da kyau ko kaɗan.

Amma tun da shi ne wanzami ɗaya tak na ƙauyen, ba makawa sai an bar shi
ya aske ya kuma yi wa mutum aski. Ban da haka kuma mutum ne mai rowa
kuma marar kirki.

%%K t07-c1-08-1 p
8. --- Da sa’a, sauran \VUgras{maƙwabta} ba su yi kama da shi ba. Alal
misali \VUgras{mai yin karusa} \textsuperscript{(9)} mutum ne mai kirki,
mai himma kuma mai taimako. Jin daɗi ne mutum ya sa shi ya gyara masa
abin hawa. Aikinsa kullum yana ƙarewa a kan lokaci.

%%K t07-c1-09-1 p
9. --- Tsoho Iakobus bai yi gunaguni a kan \VUgras{mai yin sirdi}
\textsuperscript{(10)} ba ma, wanda yakan taimake shi ya ɗinka
\VUgras{kayan doki} idan aikin ya yi yawa.

%%K t07-c1-09-2 p
Mahaifina ya tambaye shi game da iyalinsa : « Kowa yana lafiya.
Jikanyata tana \VUgras{makaranta} \textsuperscript{(12)}.
\VUgras{Malamarta} \textsuperscript{(13)} ta gamsu da ita ƙwarai,
\VUgras{ɗaliba} \textsuperscript{(14)} ce mai kyau. Ba ta yi kama da
ɗana mugun yaro ba ; wasa kaɗai yake tunani. \VUgras{Malami}
\textsuperscript{(15)} bai iya sa shi ya koyi komai ba. »

%%K t07-tit-4 sub
\VUsaut{3.0mm}
\VUpk{10.2pt}{40}{Hirar ƙauye.}
\VUsaut{3.0mm}

%%K t07-c1-10-1 p
10. --- Sa’an nan tsohon ya ba mu labaran ƙauyen. Za a gina sabuwar
makaranta, domin wadda take a majalisar gari ta zama ƙarama ƙwarai.
Ban da haka wannan abu mai sauƙi ne, domin siyan ɓangaren lambun mai
niƙa zai isa. Wannan zai amfani mutumin nan mai tausayi ƙwarai. Saboda
sanyi, \VUgras{duwatsun niƙa} na \VUgras{gidan niƙa}
\textsuperscript{(16)} ba sa iya niƙa alkama kuma, domin
\VUgras{keke} \textsuperscript{(17)} ya tsaya cak saboda \VUgras{ƙanƙara}
\textsuperscript{(25)} ta \VUgras{ƙaramin kogi} \textsuperscript{(23)}.

%%K t07-c1-11-1 p
11. --- « Game da gine-gine, shin kun ga sabuwar \VUgras{cocinmu}
\textsuperscript{(18)} ? Tana da ban sha’awa ƙwarai a ƙarƙashin
mayafinta na dusar ƙanƙara. \VUgras{Hankaka} \textsuperscript{(19)},
waɗanda ba su gane tsohuwar hasumiyarsu kuma ba, suna zama cikin baƙin
ciki a kan babbar bishiyar \VUgras{makabarta} \textsuperscript{(20)}. »

%%K t07-c1-12-1 p
12. --- Wannan tunani game da makabarta ya tunatar da mai gyaran
takalma mutanen da mahaifina ya sani kuma suka mutu bara. « Kai,
\VUgras{ƙaburbura} \textsuperscript{(21)} nawa ne, malamina, aka ƙara
kan sauran, a ƙarƙashin \VUgras{dogon icen makabarta}
\textsuperscript{(22)} ! Mutuwa ta yanke tsohon magatakardarmu. Dukan
mutanen ƙauye sun halarci \VUgras{jana’izarsa.} »

%%K t07-c1-13-1 p
13. --- « Ga wanda ya gaje shi yana zamiya a kan ƙaramin kogin. Ya zo
yanzu ya sa na gyara masa igiyar \VUgras{takalman ƙanƙararsa}
\textsuperscript{(26)} wadda ta karye. »

%%K t07-c1-14-1 p
14. --- « Ga kuma a gaɓar ƙaramin kogin, sabon \VUgras{mai gadin
ƙauye} \textsuperscript{(27)}. Tsohon ɗan sanda ne wanda yake tsoratar da
ƴan iskan ƙauyen. Suna gudu nan take idan suka hangi \VUgras{murfin
ƙafarsa} \textsuperscript{(29)} da \VUgras{hular sojansa}
\textsuperscript{(28)}. »

%%K t07-c1-15-1 p
15. --- « Kai ! ga tsoho Iosef yana tuƙa \VUgras{ƙaramar karusar icen
wuta} \textsuperscript{(30)} domin mai gasa burodi. --- Gaskiya ne, in
ji mahaifina, na gane jarginsa na \VUgras{alfadarai}
\textsuperscript{(31)}. Ina buƙatar yin magana da shi, zan yi amfani da
damar. Sai mun sake gani, baba Iakobus. »

%%K t07-c1-16-1 p
16. --- Daidai lokacin da muka fita, yaran ƙauyen suka bar makaranta
suka gaggauta zuwa \VUgras{wurin zamiya} \textsuperscript{(33)} cikin
gudu, cikin haɗarin faɗuwa da karya gaɓa a \VUgras{faɗuwa}
\textsuperscript{(34)}. Ɗaya daga cikinsu ya ɗauko \VUgras{karusar
ƙanƙararsa} \textsuperscript{(32)} wurin mai yin karusa, ya zaunar da
ɗaya daga cikin abokansa a cikinta, ya tafi a guje.

%%K t07-c1-17-1 p
17. --- Waɗansu suka gaggauta zuwa \VUgras{tudun dusar ƙanƙara}
\textsuperscript{(37)} kusa da \VUgras{gada} \textsuperscript{(40)} suka
fara jifan \VUgras{mutum-mutumin dusar ƙanƙara}
\textsuperscript{(39)} da suka kafa jiya kuma da suka sa masa hula da
bututun taba da jakar makaranta a kafaɗa.

%%K t07-c1-18-1 p
18. --- Ba su damu sosai da iska mai sanyi da sanyin ba. Yawancinsu ba
su ma da \VUgras{mayafin wuya} \textsuperscript{(35)}, kuma, domin su
sake ɗumama bayan yaƙin \VUgras{ƙwallon dusar ƙanƙara}
\textsuperscript{(38)}, tafin hannu na \VUgras{kastanya}
\textsuperscript{(42)} mai zafi ƙwarai da aka saya wurin tsohuwar
\VUgras{mai sayarwa} Jakobina ya ishe su, ita wadda take rawar jiki
saboda sanyi a ƙarƙashin \VUgras{ƙaramin mayafinta}
\textsuperscript{(43)} mai sirara ƙwarai.

%%K t07-tit-5 sub
\VUsaut{3.0mm}
\VUcentre{12.0pt}{0}{\textit{Fage na Biyu.}}
\VUsaut{3.0mm}

%%K t07-tit-6 sub
\VUpk{10.2pt}{40}{Gidan gona a lokacin bazara.}
\VUsaut{4.0mm}

%%K t07-c2-01-1 p
1. --- Duk da dukan jin daɗin hunturu, muna gaisuwar dawowar
\VUgras{bazara} da farin ciki mai zurfi.

%%K t07-c2-01-2 p
Kai, hoto mai faranta rai ne farfajiyar gona take yi, da yamma, a
watan Mayu !

%%K t07-c2-02-1 p
2. --- A sararin sama \VUgras{rana} \textsuperscript{(1)} tana faɗuwa a
hankali, tana cika \VUgras{ƙauye} \textsuperscript{(3)} da \VUgras{haskenta}
\textsuperscript{(2)} mai launin jan shunayya. \VUgras{Manomi mai
himma} \textsuperscript{(6)} yana gaggawar nomawa \VUgras{gonarsa}
\textsuperscript{(4)}.

%%K t07-c2-02-2 p
Da \VUgras{garmarsa} \textsuperscript{(5)} da aka jarge wa
\VUgras{doki} \textsuperscript{(7)} mai ƙarfi, yana zana \VUgras{layin
noma} \textsuperscript{(8)} inda \VUgras{mai shuka} \textsuperscript{(9)}
zai jefa \VUgras{iri} da cikakken tafin hannu, irin da zai ba da kan
hatsi na zinare.

%%K t07-c2-03-1 p
3. --- \VUgras{Masu lauje} \textsuperscript{(11)} da aka ba su
laujensu sun yanka ciyawar filin, masu shanya sun bushe
\VUgras{busasshiyar ciyawa} \textsuperscript{(10)} suna juya ta da
\VUgras{sandunan mai yatsa} \textsuperscript{(12)} da matsefi, ga su
kuwa suna dawowa.

%%K t07-c2-03-2 p
Waɗansu suna tafiya gaban abin hawa mai nauyi da shanu suke ja ; suna
ɗauke da kayan aikinsu a kafaɗa. Waɗansu, masu ragwanci, sun zauna
cikin nutsuwa a kan busasshiyar ciyawa mai ƙamshi.

%%K t07-c2-04-1 p
4. --- Suna wucewa suna yi wa \VUgras{mai lambu} \textsuperscript{(16)}
gaisuwar abokantaka, shi da yake aiki a \VUgras{gonar ƴaƴan itace}
\textsuperscript{(13)} yana sassaƙa \VUgras{itatuwan ƴaƴan itace} yana
kuma yin gyaran \VUgras{itatuwan fir} \textsuperscript{(14)}.
\VUgras{Itatuwan furam} \textsuperscript{(15)} suna fara yin fure, kuma,
a \VUgras{gonar ganyayyaki} \textsuperscript{(17)} da aka yi mata taki
da kyau, ganyayyaki da kabeji da salati sun riga sun yi kyau.

%%K t07-c2-04-2 p
A kan ƙaramin bango, \VUgras{dawisu} \textsuperscript{(18)} mai kyau
yana buɗe wutsiyarsa, domin a yaba wa gashinsa masu kyau ƙwarai.

%%K t07-tit-7 sub
\VUpk{10.2pt}{40}{Farfajiyar gona.}
\VUsaut{3.0mm}

%%K t07-c2-05-1 p
5. --- Ƙofofin \VUgras{gidan dawaki} \textsuperscript{(19)} suna a buɗe
kuma ana ganin mariƙin ciyawa da \VUgras{shimfiɗar}
\textsuperscript{(23)} \VUgras{godiya} \textsuperscript{(21)} wadda
\VUgras{mai kula da doki} \textsuperscript{(20)} ya kwance yanzu.
\VUgras{Ɗan doki} \textsuperscript{(22)} mai ban dariya ƙwarai da
doguwar ƙafafunsa, yana bin mahaifiyarsa yana tsalle.

%%K t07-c2-06-1 p
6. --- Kusa da \VUgras{rijiya} \textsuperscript{(29)}, \VUgras{shanu}
\textsuperscript{(25)} suna zuwa su sha a \VUgras{kwatarniyar}
\textsuperscript{(26)} katako wadda take musu mashayi. \VUgras{Ɗan sa}
\textsuperscript{(24)} yana amfani da damar ya sha nono, \VUgras{bijimi}
\textsuperscript{(27)} kuwa yana jin ƙamshin ciyawa mai kyau, yana yin
ruri da farin ciki yana kuma ɗaga kansa mai \VUgras{ƙahoni}
\textsuperscript{(28)} masu ƙarfi cikin alfahari.

%%K t07-c2-07-1 p
7. --- Kowa yana aiki. \VUgras{Mai aikin gida} \textsuperscript{(32)}
tana riƙe da \VUgras{igiyar} \textsuperscript{(31)} rijiya da hannu.
Tana ɗiban ruwa da \VUgras{guga} \textsuperscript{(30)} domin ta shayar
da dabbobi. Kusa da \VUgras{rumfar hatsi} \textsuperscript{(33)},
wani mutum yana share ciyawa, a gaban ƙofar \VUgras{gida}
\textsuperscript{(34)} kuwa tsohuwar \VUgras{mai kaɗi}
\textsuperscript{(35)} tana kaɗi da \VUgras{keken kaɗinta} da
\VUgras{sandarta,} tana kaɗa \VUgras{lilin} ko \VUgras{rama} kamar yadda
ake yi a zamanin dā.

%%K t07-tit-8 sub
\VUsaut{3.0mm}
\VUpk{10.2pt}{40}{Manomi.}
\VUsaut{3.0mm}

%%K t07-c2-08-1 p
8. --- \VUgras{Manomi} \textsuperscript{(36)} yana jagorantar dukan
waɗannan ayyuka yana kuma ba wa ma’aikatansa umarni. Yana ba da
shawara a saka \VUgras{matsefin gona} \textsuperscript{(40)} da
\VUgras{fartanya} \textsuperscript{(39)} a ƙarƙashin rufi, waɗanda aka
manta da su kusa da \VUgras{ɗakin} \textsuperscript{(38)}
\VUgras{kare} \textsuperscript{(37)}.

%%K t07-c2-09-1 p
9. --- Kiransa yana tsoratar da \VUgras{tantabara}
\textsuperscript{(41)}, wadda take shiga \VUgras{gidan tantabaru}
\textsuperscript{(42)} da dukan fikafikanta, da \VUgras{kaji}
\textsuperscript{(43)} waɗanda suke gaggawar komawa \VUgras{gidan kaji}
\textsuperscript{(44)}.

%%K t07-c2-10-1 p
10. --- A gaban \VUgras{gidan tumaki} \textsuperscript{(48)}, ga
tsohon \VUgras{makiyayi} \textsuperscript{(45)} wanda yake mayar da
\VUgras{garkensa} \textsuperscript{(49)}. Yana riƙe da \VUgras{sandarsa}
\textsuperscript{(46)}, wadda ba ta taɓa rasa masa ba, kamar
\VUgras{rigar samansa} \textsuperscript{(47)} da ta koma ba ta da
launi, yana kai wa turke \VUgras{raguna} \textsuperscript{(50)}, da
\VUgras{tumaki} \textsuperscript{(53)}, da \VUgras{ƴan tumaki}
\textsuperscript{(54)}, da \VUgras{awaki} \textsuperscript{(51)} da
\VUgras{ƴan awaki} \textsuperscript{(52)}. \VUgras{Karensa}
\textsuperscript{(55)} mai aminci Argus ya zo a ƙarshe yana
ƙarfafa masu jinkiri da haushinsa.

%%K t07-c2-11-1 p
11. --- Dukan wannan hayaniya da motsi ba sa dagula natsuwar
\VUgras{saniya mai nono} \textsuperscript{(56)} mai kyau wadda take kaɗa
\VUgras{ƙararrawarta} \textsuperscript{(57)}, alhali ana tatsarta a
gaban ƙofar \VUgras{turke} \textsuperscript{(58)}.

%%K t07-c2-12-1 p
12. Ta yi kama da wadda ta gane cewa dole ta ba da nononta, domin
\VUgras{matar manomi} \textsuperscript{(60)} ta iya shirya \VUgras{man
shanu} a cikin \VUgras{abin kaɗa man shanu} \textsuperscript{(61)} ko
kyawawan \VUgras{cuku} \textsuperscript{(64)}, waɗanda suke ɗiga daga
kantar da aka ajiye a kan \VUgras{babban kwano}
\textsuperscript{(63)}.

%%K t07-c2-13-1 p
13. --- Ana kula da dukan \VUgras{wurin nono} \textsuperscript{(59)}
cikin tsabta ƙwarai ta wurin \VUgras{ma’aikacin gona}
\textsuperscript{(65)}, shi wanda ya wanke \VUgras{tuluna na nono} \textsuperscript{(62)} yanzu
a cikin babbar guga wadda yake ɗauke da ita a hannu.

%%K t07-tit-9 sub
\VUsaut{3.0mm}
\VUpk{10.2pt}{40}{Gidan tsuntsaye.}
\VUsaut{3.0mm}

%%K t07-c2-14-1 p
14. --- Da takalmansa masu kauri na katako, yana tafiya da ɗan nauyi,
haka kuwa yake sa \VUgras{talotalo} \textsuperscript{(68)}, da
\VUgras{agwagwa mata} \textsuperscript{(66)}, da \VUgras{agwagwa maza}
\textsuperscript{(67)} da \VUgras{manyan agwagwa} \textsuperscript{(69)}
su gudu, waɗanda suke buɗe \VUgras{bakunansu} \textsuperscript{(70)}
sosai suna kuka da fargaba.

%%K t07-c2-14-2 p
\VUgras{Zakara} \textsuperscript{(71)}, mai son faɗa, yana ɗaga
\VUgras{jan tsefensa} \textsuperscript{(72)}, yana girgiza gashin
\VUgras{wutsiyarsa} \textsuperscript{(73)} yana kuma yin « kukuruku »
mai ƙara.

%%K t07-c2-15-1 p
15. --- \VUgras{Kaza mai ƴaƴa} \textsuperscript{(74)} tana yawo a kan
\VUgras{taki} \textsuperscript{(81)} tare da \VUgras{ƴaƴanta}
\textsuperscript{(75)}. \VUgras{Ɗan alade} \textsuperscript{(78)} yana
gudu zuwa wurin mahaifiyarsa, \VUgras{aladiya}
\textsuperscript{(77)} babba wadda muke gani kusa da gidan aladu.

%%K t07-c2-16-1 p
16. --- A cikin kejinsu, \VUgras{zomaye} \textsuperscript{(79)} suna
cin \VUgras{ciyawa} \textsuperscript{(80)} mai laushi da ƴar manomi take
kawo musu da haɗama. Ioanina tana son zomayenta ƙwarai kuma, kowace
rana, bayan ta ɗauko sabbin \VUgras{ƙwai} \textsuperscript{(76)} a gidan
kaji, takan zo ta ziyarci ƙananan abokanta.

\VUsaut{6.0mm}
\VUfilet{22mm}
